\chapter{Python for Data Science: the Core Language}
\companion{Corey Schafer: Python Tutorial for Beginners (search the series)}{\yts{Corey+Schafer+Python+Tutorial+for+Beginners+Strings}}

The pre-placement talk said programming MCQs in the online test ``will be quite easy'', giving the example of converting a string to an int and
spotting the right output. Round 3 then puts you in a live notebook. This chapter teaches the language features that such questions test,
assuming no prior Python. Every output shown here was checked by running the code (Python 3.11).

\section{Values, types and variables}
A \textbf{variable} is a name that refers to an object. Every object has a \textbf{type}:
\begin{center}\small
\begin{tabularx}{\linewidth}{l l X}
\toprule
Type & Example & Notes\\
\midrule
\code{int} & \code{42}, \code{-7} & unlimited size\\
\code{float} & \code{3.14}, \code{1e-3} & 64-bit; \code{0.1 + 0.2 == 0.3} is \code{False}\\
\code{bool} & \code{True}, \code{False} & a subclass of int: \code{True + True == 2}\\
\code{str} & \code{"Pune"} & immutable sequence of characters\\
\code{list} & \code{[1, "a", 2.5]} & ordered, mutable\\
\code{tuple} & \code{(1, 2)}, \code{(1,)} & ordered, immutable; \code{(1)} is just \code{1}\\
\code{dict} & \code{\{"city": "Pune"\}} & key $\to$ value; insertion-ordered\\
\code{set} & \code{\{1, 2, 3\}} & unordered, unique elements\\
\code{NoneType} & \code{None} & ``no value''\\
\bottomrule
\end{tabularx}
\end{center}

\subsection{Converting between types (the OA's favourite)}
\begin{lstlisting}
int("12") + 1        # 13
int(3.9), int(-3.9)  # (3, -3)   truncates toward zero
int("3.5")           # ValueError: invalid literal for int() with base 10: '3.5'
int(float("3.5"))    # 3
"3" + "4"            # '34'      string concatenation
"3" * 3              # '333'
str(10) + "5"        # '105'
int("101", 2)        # 5         binary
round(2.5), round(3.5)   # (2, 4)  banker's rounding: ties go to the even number
round(2.675, 2)      # 2.67      2.675 is stored as 2.67499999...
\end{lstlisting}

\subsection{Truthiness}
In an \code{if}, these are false: \code{0}, \code{0.0}, \code{""}, \code{[]}, \code{\{\}}, \code{set()}, \code{None}, \code{False}. Everything else is true,
including \code{"0"} and \code{[0]}. So \code{bool("0")} is \code{True}.

\section{Operators}
\begin{lstlisting}
7 // 2, -7 // 2      # (3, -4)   floor division rounds DOWN, not toward zero
7 % 3, -7 % 3        # (1, 2)    result takes the sign of the divisor
5 / 2, 4 / 2         # (2.5, 2.0)  / always returns a float
2 ** 3 ** 2          # 512       ** is right-associative: 2 ** 9
divmod(17, 5)        # (3, 2)
1 == 1.0             # True      equal values
[] == [], [] is []   # (True, False)  == compares values, `is` compares identity
\end{lstlisting}

\section{Strings}
Strings are immutable sequences: indexing and slicing work like lists, but you cannot assign to a character.
\begin{lstlisting}
s = "DataScience"
s[0], s[-1]          # ('D', 'e')
s[2:6]               # 'taSc'      start inclusive, stop exclusive
s[-3:]               # 'nce'
s[::2]               # 'DtSine'    every second character
s[::-1]              # 'ecneicSataD'  reversed
len(s)               # 11
"a,b,,c".split(",")  # ['a', 'b', '', 'c']
" x ".strip()        # 'x'
"-".join(["a", "b"]) # 'a-b'
"abc".find("z")      # -1   (index() would raise ValueError)
"12".isdigit(), "1.2".isdigit()   # (True, False)
f"{3.14159:.2f}"     # '3.14'
s[0] = "d"           # TypeError: 'str' object does not support item assignment
\end{lstlisting}

\section{Lists, tuples, dicts and sets}
\begin{lstlisting}
L = [3, 1, 2]
L.sort()             # sorts in place and RETURNS None
sorted(L, reverse=True)   # [3, 2, 1]  returns a new list
L.append(4); L.extend([5, 6]); L.pop(); L.insert(0, 9)
nums = [1, 2, 3, 4]
nums[1:-1], nums[::-1]    # ([2, 3], [4, 3, 2, 1])
a, *rest = [1, 2, 3, 4]   # a = 1, rest = [2, 3, 4]

d = {"a": 1}
d.get("b", 0)        # 0   (d["b"] would raise KeyError)
for k, v in d.items(): ...
{k: v for k, v in [("a", 1), ("b", 2)] if v > 1}   # {'b': 2}

s, s2 = {1, 2, 3}, {2, 3, 4}
s & s2, s | s2, s - s2, s ^ s2   # ({2, 3}, {1, 2, 3, 4}, {1}, {1, 4})
\end{lstlisting}
Membership: \code{x in list} is $O(n)$; \code{x in set} and \code{x in dict} (keys) are $O(1)$ on average because they use hashing. Dict keys and set elements
must be hashable (immutable): a tuple can be a key, a list cannot.

\section{Mutability, references and copying}
This is the source of most ``what is the output'' traps.
\begin{lstlisting}
a = [1, 2, 3]; b = a; b.append(4); print(a)     # [1, 2, 3, 4]  same object
a = [1, 2, 3]; b = a[:]; b.append(4); print(a)  # [1, 2, 3]     a copy
m = [[0] * 2] * 2; m[0][0] = 5; print(m)        # [[5, 0], [5, 0]]  both rows are ONE list
m = [[0] * 2 for _ in range(2)]                 # correct: two separate rows

import copy
a = [[1, 2], [3]]
b = copy.copy(a)        # shallow: new outer list, same inner lists
c = copy.deepcopy(a)    # deep: everything copied
a[0].append(9)
print(b, c)             # [[1, 2, 9], [3]] [[1, 2], [3]]

t = ([1], 2); t[0].append(3); print(t)   # ([1, 3], 2)  tuple fixed, its list is mutable
\end{lstlisting}
\textbf{Function arguments} are passed by object reference (``call by sharing''): a function can mutate a mutable argument, but rebinding the name inside does not
affect the caller.
\begin{lstlisting}
x = [1, 2, 3]
def mod(l):  l += [4]       # in-place: caller sees it
def mod2(l): l = l + [5]    # new list bound to local name: caller does not
mod(x); mod2(x); print(x)   # [1, 2, 3, 4]
\end{lstlisting}

\subsection{The mutable default argument trap}
\begin{lstlisting}
def f(x, L=[]):
    L.append(x)
    return L
print(f(1), f(2))     # [1, 2] [1, 2]   the default list is created ONCE
def f(x, L=None):     # fix
    if L is None: L = []
    L.append(x); return L
\end{lstlisting}

\section{Control flow and comprehensions}
\begin{lstlisting}
[i * i for i in range(5) if i % 2 == 0]       # [0, 4, 16]
{w: len(w) for w in ["ml", "python"]}          # {'ml': 2, 'python': 6}
list(enumerate("ab", start=1))                  # [(1, 'a'), (2, 'b')]
list(zip([1, 2, 3], "ab"))                      # [(1, 'a'), (2, 'b')]   stops at shortest
any([0, 0, 1]), all([1, 1, 0]), all([])        # (True, False, True)
\end{lstlisting}

\section{Functions: \code{*args}, \code{**kwargs}, lambda, scope}
\begin{lstlisting}
def k(*args, **kw): return args, kw
k(1, 2, a=3)                 # ((1, 2), {'a': 3})
list(map(lambda v: v * 2, [1, 2, 3]))        # [2, 4, 6]
list(filter(None, [0, 1, "", 2]))            # [1, 2]
sorted(["b", "A", "c"])                      # ['A', 'b', 'c']   uppercase sorts first
sorted(["b", "A", "c"], key=str.lower)       # ['A', 'b', 'c']
sorted({"b": 2, "a": 3}.items(), key=lambda kv: kv[1], reverse=True)  # [('a', 3), ('b', 2)]
\end{lstlisting}
\textbf{Scope (LEGB)}: names are looked up in Local, Enclosing function, Global, Built-in order. Assigning to a global inside a function needs \code{global}.
\textbf{Late binding of closures}:
\begin{lstlisting}
funcs = [lambda: i for i in range(3)]; [f() for f in funcs]        # [2, 2, 2]
funcs = [lambda i=i: i for i in range(3)]; [f() for f in funcs]    # [0, 1, 2]
\end{lstlisting}

\section{Iterators and generators}
A \textbf{generator} produces values lazily, one at a time, with \code{yield}; it remembers where it paused. It uses constant memory, ideal for streaming a
10 GB log file line by line.
\begin{lstlisting}
def countdown(n):
    while n > 0:
        yield n
        n -= 1
list(countdown(3))         # [3, 2, 1]
g = (i for i in range(3))  # generator expression
sum(g), sum(g)             # (3, 0)   a generator is exhausted after one pass
\end{lstlisting}
\code{next(g)} on an exhausted generator raises \code{StopIteration}.

\section{Decorators}
A decorator is a function that takes a function and returns a new function that adds behaviour (logging, timing, caching, retries, auth checks).
\begin{lstlisting}
import functools, time
def timer(fn):
    @functools.wraps(fn)                 # keep fn's name and docstring
    def wrapper(*a, **k):
        t = time.perf_counter()
        result = fn(*a, **k)
        print(f"{fn.__name__} took {time.perf_counter() - t:.4f}s")
        return result
    return wrapper

@timer                 # same as: add = timer(add)
def add(a, b): return a + b
add(2, 3)              # prints the timing line, returns 5
\end{lstlisting}
\code{@functools.lru\_cache} is a built-in caching decorator.

\section{Exceptions}
\begin{lstlisting}
try:
    value = int(row["salary"])
except (ValueError, TypeError):
    value = None
else:
    pass                 # runs if no exception
finally:
    pass                 # always runs
\end{lstlisting}
Catch specific exceptions, not bare \code{except:}.

\section{Classes in one page}
\begin{lstlisting}
class Candidate:
    count = 0                              # class attribute, shared
    def __init__(self, name, years):
        self.name, self.years = name, years  # instance attributes
        Candidate.count += 1
    def is_senior(self):
        return self.years >= 5
    def __repr__(self):
        return f"Candidate({self.name!r}, {self.years})"
\end{lstlisting}
\code{\_\_eq\_\_} defines \code{==}; without it, \code{==} falls back to identity.

\section{Performance notes}
Vectorise with NumPy/pandas instead of Python loops (Chapter 25). Use sets/dicts for membership. Use generators for big streams. The \textbf{GIL} (global interpreter
lock) lets only one thread execute Python bytecode at a time, so CPU-bound work is parallelised with \code{multiprocessing} (or done in NumPy, which releases the
GIL), while threads suit I/O-bound work.

\begin{practical}[: Python in the InfoEdge rounds]
OA: short output-prediction MCQs on types, slicing, conversion, mutability. Round 3: data handling in pandas and an sklearn model (Chapters 26--27); clean,
readable code with small functions and clear names counts.
\end{practical}

\stuck{Ned Batchelder: Facts and Myths about Python Names and Values}{\yts{Ned+Batchelder+Facts+and+Myths+about+Python+Names+and+Values}}
\section{Interview and OA questions}

\begin{qa}{\pyq{2026--27 talk: string to int} What do \code{int("12") + 1}, \code{int("3.5")}, \code{int(3.9)} and \code{"3" + "4"} give?}
\oneline{13; a \code{ValueError}; 3; and the string \code{'34'}.}
\why \code{int()} parses a string only if it is a valid integer literal; ``3.5'' is not, so convert via \code{float} first: \code{int(float("3.5")) == 3}.
\code{int(3.9)} truncates toward zero (not rounding). \code{+} on strings concatenates.
\pushes \emph{``How do you safely convert a messy salary column?''} \code{pd.to\_numeric(df["salary"], errors="coerce")} turns unparsable values into NaN.
\end{qa}

\begin{qa}{What is printed: \code{a = [1,2,3]; b = a; b.append(4); print(a)}? And with \code{b = a[:]}?}
\oneline{\code{[1, 2, 3, 4]} in the first case (both names refer to the same list) and \code{[1, 2, 3]} in the second (slicing makes a shallow copy).}
\why Assignment never copies in Python; it binds another name to the same object. Check with \code{a is b}.
\end{qa}

\begin{qa}{What does \code{m = [[0]*2]*2; m[0][0] = 5; print(m)} print, and how do you fix it?}
\oneline{\code{[[5, 0], [5, 0]]}, because the outer \code{*2} repeats a reference to the same inner list; build rows with a comprehension: \code{[[0]*2 for \_ in range(2)]}.}
\end{qa}

\begin{qa}{Explain shallow vs deep copy with an example.}
\oneline{A shallow copy creates a new outer container that still references the same inner objects; a deep copy recursively copies everything.}
\worked After \code{a[0].append(9)}: shallow copy shows \code{[[1, 2, 9], [3]]}, deep copy \code{[[1, 2], [3]]}. In pandas, \code{df.copy()} is deep by default for data.
\end{qa}

\begin{qa}{What is wrong with \code{def f(x, L=[])}? Show the output of \code{f(1), f(2)}.}
\oneline{The default list is created once at definition time and shared across calls, so both calls append to the same list: \code{[1, 2] [1, 2]}; use \code{L=None} and create
a new list inside.}
\end{qa}

\begin{qa}{Lists vs tuples vs sets vs dicts: when do you use each?}
\oneline{List for an ordered, changeable sequence; tuple for a fixed record or a dictionary key; set for uniqueness and fast membership; dict for key-to-value lookups.}
\why Membership in a list is $O(n)$, in a set or dict $O(1)$ average. Tuples are hashable (if their contents are), so \code{(city, title)} can key a dict of average salaries.
\end{qa}

\begin{qa}{Give the outputs: \code{7 // 2}, \code{-7 // 2}, \code{-7 \% 3}, \code{2 ** 3 ** 2}, \code{5 / 2}, \code{round(2.5)}.}
\oneline{3, $-4$, 2, 512, 2.5, 2.}
\why Floor division rounds toward $-\infty$; \code{\%} has the divisor's sign so that \code{(a // b) * b + a \% b == a}; \code{**} is right-associative; \code{/} is true division;
\code{round} uses banker's rounding.
\end{qa}

\begin{qa}{Slice \code{s = "DataScience"}: \code{s[2:6]}, \code{s[-3:]}, \code{s[::-1]}, \code{s[::2]}.}
\oneline{\code{'taSc'}, \code{'nce'}, \code{'ecneicSataD'}, \code{'DtSine'}.}
\end{qa}

\begin{qa}{What is the difference between \code{==} and \code{is}?}
\oneline{\code{==} compares values (calls \code{\_\_eq\_\_}); \code{is} checks whether two names refer to the same object; use \code{is} only for singletons like \code{None}.}
\worked \code{[] == []} is \code{True}; \code{[] is []} is \code{False}.
\end{qa}

\begin{qa}{What is a generator, and why would you use one on a 10 GB log of job applications?}
\oneline{A function with \code{yield} that produces values lazily one at a time and keeps its state between calls; it lets you process the file line by line in constant memory.}
\worked
\begin{lstlisting}
def parse(path):
    with open(path) as f:
        for line in f:                    # files are iterators too
            uid, job, ts = line.rstrip("\n").split("\t")
            yield uid, job
applies_per_job = Counter(job for _, job in parse("applies.tsv"))
\end{lstlisting}
\end{qa}

\begin{qa}{What is a decorator? Write one that times a function.}
\oneline{A function that wraps another function to add behaviour without changing its code; \code{@timer} above is equivalent to \code{add = timer(add)}.}
\why Use \code{functools.wraps} so the wrapped function keeps its name. Real uses: caching (\code{lru\_cache}), retries for flaky API calls, logging model latency.
\end{qa}

\begin{qa}{\deep Why does \code{[lambda: i for i in range(3)]} return functions that all give 2?}
\oneline{Closures capture variables, not values; all lambdas refer to the same \code{i}, which is 2 after the loop ends; bind the current value with a default argument
\code{lambda i=i: i}.}
\end{qa}

\begin{qa}{What are \code{*args} and \code{**kwargs}?}
\oneline{\code{*args} collects extra positional arguments into a tuple, \code{**kwargs} extra keyword arguments into a dict; also used to unpack when calling,
\code{f(*lst, **dct)}.}
\worked \code{k(1, 2, a=3)} returns \code{((1, 2), \{'a': 3\})}.
\end{qa}

\begin{qa}{Count word frequencies in a list of job titles without pandas.}
\oneline{Use \code{collections.Counter} on the split words.}
\worked
\begin{lstlisting}
from collections import Counter
titles = ["Data Scientist", "Senior Data Scientist", "Data Analyst"]
c = Counter(w.lower() for t in titles for w in t.split())
c.most_common(2)        # [('data', 3), ('scientist', 2)]
\end{lstlisting}
\end{qa}

\begin{qa}{What does the GIL mean for data science code?}
\oneline{Only one thread runs Python bytecode at a time, so pure-Python CPU-bound loops do not speed up with threads; use vectorised NumPy/pandas (which run in C and
release the GIL), \code{multiprocessing}, or libraries with their own parallelism (\code{n\_jobs=-1}).}
\end{qa}

\begin{qa}{\deep Why is \code{0.1 + 0.2 == 0.3} false, and how should you compare floats?}
\oneline{Because 0.1 and 0.2 have no exact binary representation, so their sum is $0.30000000000000004$; compare with a tolerance using \code{math.isclose} or
\code{np.isclose}.}
\end{qa}

\begin{qa}{What does \code{L.sort()} return? What about \code{sorted(L)}?}
\oneline{\code{L.sort()} sorts in place and returns \code{None}; \code{sorted(L)} returns a new sorted list and leaves \code{L} unchanged.}
\worked \code{print(L.sort())} prints \code{None} (a classic OA trap).
\end{qa}

\begin{qa}{Sort a dict of skill counts by count, descending.}
\oneline{\code{sorted(d.items(), key=lambda kv: kv[1], reverse=True)}.}
\end{qa}

\begin{qa}{What happens with \code{d["missing"]} vs \code{d.get("missing")}?}
\oneline{The first raises \code{KeyError}; \code{get} returns \code{None} or a given default; for counting, \code{defaultdict(int)} or \code{Counter} avoid the issue.}
\end{qa}

\begin{qa}{\deep A function receives a list and does \code{l = l + [5]}; another does \code{l += [5]}. Which changes the caller's list and why?}
\oneline{Only \code{l += [5]}, because for lists \code{+=} mutates in place (\code{list.\_\_iadd\_\_}), while \code{l + [5]} builds a new list and rebinds the local name.}
\end{qa}
